
\documentclass[aspectratio=169,11pt]{beamer}

\usetheme{AnnArbor}
\usepackage{fontspec}
\usepackage[vietnamese]{babel}
\usepackage{amsmath,amssymb,booktabs,array,multicol}
\usepackage{listings}
\usepackage{xcolor}
\usepackage{graphicx}
\usepackage{hyperref}

\definecolor{CodeBg}{RGB}{248,248,248}
\definecolor{TitleBlue}{RGB}{38,79,129}
\definecolor{AccentGreen}{RGB}{35,125,93}
\definecolor{SoftGray}{RGB}{80,80,80}
\definecolor{QuizOrange}{RGB}{180,95,30}

\setbeamercolor{title}{fg=TitleBlue}
\setbeamercolor{frametitle}{fg=TitleBlue}
\setbeamercolor{structure}{fg=AccentGreen}
\setbeamercolor{block title}{fg=white,bg=TitleBlue}
\setbeamercolor{block body}{bg=blue!3}
\setbeamercolor{alertblock title}{fg=white,bg=QuizOrange}
\setbeamercolor{alertblock body}{bg=orange!6}
\setbeamertemplate{navigation symbols}{}
\setbeamertemplate{footline}[frame number]
\setbeamertemplate{itemize item}{\raise0.15ex\hbox{\tiny$\blacktriangleright$}}
\setbeamertemplate{itemize subitem}{\raise0.15ex\hbox{\tiny$\bullet$}}

\lstdefinestyle{py}{
  language=Python,
  basicstyle=\ttfamily\scriptsize,
  backgroundcolor=\color{CodeBg},
  frame=single,
  framerule=0.2pt,
  rulecolor=\color{gray!35},
  breaklines=true,
  showstringspaces=false,
  keywordstyle=\color{TitleBlue}\bfseries,
  commentstyle=\color{SoftGray}\itshape,
  stringstyle=\color{AccentGreen}
}

\newcommand{\key}[1]{\textbf{\textcolor{AccentGreen}{#1}}}
\newcommand{\smallnote}[1]{\vspace{0.25em}{\footnotesize\color{SoftGray}#1}}
\newcommand{\quiztitle}[1]{\textbf{\textcolor{QuizOrange}{#1}}}
\newcommand{\imageplaceholder}[2]{%
\begin{center}
  \fcolorbox{SoftGray}{gray!5}{%
  \begin{minipage}[c][0.55\textheight][c]{0.84\linewidth}
  \centering
  {\Large\textcolor{TitleBlue}{Image Placeholder}}\\\[1.0em]
  {\normalsize\textbf{#2}}\\\[1.2em]
  {\footnotesize Source path in Markdown:}\\\[0.35em]
  {\small\texttt{\detokenize{#1}}}
  \end{minipage}}
\end{center}
}

\title[Giới thiệu NumPy]{Giới thiệu NumPy}
\subtitle{Mảng, vector hóa, broadcasting và tính toán số}
\author{Duc-Minh Vu}
\institute{SLSCM Lab -- FDA -- NEU}
\date{\today}

\begin{document}

\begin{frame}
  \titlepage
\end{frame}

\begin{frame}{Lộ trình bài học}
\begin{columns}[T]
\column{0.48\linewidth}
\begin{enumerate}
  \item NumPy là gì
  \item Vì sao NumPy hữu ích
  \item Cài đặt và import
  \item Tạo mảng
  \item Thuộc tính mảng
  \item Indexing và slicing
  \item Reshaping, resizing, stacking, splitting
\end{enumerate}
\column{0.48\linewidth}
\begin{enumerate}
  \setcounter{enumi}{7}
  \item Broadcasting
  \item Số học và tổng hợp
  \item Universal functions
  \item Toán tử Boolean
  \item Linear algebra
  \item Số ngẫu nhiên và thống kê
  \item Tích hợp và lỗi thường gặp
\end{enumerate}
\end{columns}
\end{frame}

\begin{frame}{Mục tiêu học tập}
\small
Sau bài học này, người học có thể:
\begin{itemize}
  \item Giải thích vai trò của NumPy trong tính toán số.
  \item Phân biệt mảng NumPy với list Python.
  \item Tạo mảng một chiều và mảng nhiều chiều.
  \item Kiểm tra \texttt{shape}, \texttt{size}, \texttt{ndim} và \texttt{dtype}.
  \item Sử dụng indexing, slicing, reshaping, stacking và splitting.
  \item Áp dụng số học vector hóa, broadcasting, tổng hợp và universal functions.
  \item Thực hiện đại số tuyến tính cơ bản và thống kê mô tả.
  \item Sinh giá trị ngẫu nhiên và hiểu tính tái lập kết quả.
  \item Giải thích cách NumPy kết nối với Pandas, SciPy và Scikit-learn.
\end{itemize}
\end{frame}

\begin{frame}[fragile]{Điều kiện tiên quyết và cài đặt}
\begin{block}{Kiến thức nền khuyến nghị}
\begin{itemize}
  \item Biến, list, vòng lặp và hàm cơ bản trong Python.
  \item Một môi trường Python có thể chạy được: Jupyter Lưu ýbook, JupyterLab, Google Colab hoặc tương tự.
\end{itemize}
\end{block}
\begin{block}{Cài đặt và import}
\begin{lstlisting}[style=py]
pip install numpy
\end{lstlisting}
\begin{lstlisting}[style=py]
import numpy as np
print(np.__version__)
\end{lstlisting}
\end{block}
\end{frame}


\begin{frame}{Cách đọc các câu lệnh NumPy}
\small
Hầu hết câu lệnh NumPy có cùng một mẫu:
\[
\texttt{np.function(input, option1=value1, option2=value2)}
\]
\begin{itemize}
  \item \texttt{np}: bí danh chuẩn của thư viện NumPy.
  \item \texttt{function}: phép toán ta muốn NumPy thực hiện.
  \item \texttt{input}: mảng, list, shape hoặc giá trị số được truyền vào hàm.
  \item \texttt{option=value}: đối số có tên dùng để điều khiển hành vi của câu lệnh.
\end{itemize}
\begin{block}{Ví dụ}
\texttt{np.arange(0, 10, 2)} tạo các giá trị bắt đầu từ 0, dừng trước 10, với bước nhảy 2.
\end{block}
\end{frame}

\begin{frame}{Mẫu câu lệnh: đối tượng, thuộc tính, phương thức, hàm}
\small
\begin{tabular}{p{0.26\linewidth}p{0.62\linewidth}}
\toprule
\textbf{Dạng} & \textbf{Ý nghĩa} \\
\midrule
\texttt{np.array([1,2,3])} & Gọi một hàm NumPy để tạo mảng. \\
\texttt{arr.shape} & Đọc một thuộc tính mô tả mảng. \\
\texttt{arr.reshape(2,3)} & Gọi một phương thức gắn với đối tượng mảng hiện có. \\
\texttt{np.mean(arr)} & Gọi một hàm NumPy để tóm tắt mảng. \\
\texttt{arr.mean()} & Phiên bản gọi theo kiểu phương thức, tương đương với nhiều phép toán phổ biến. \\
\bottomrule
\end{tabular}
\vspace{0.5em}
\begin{alertblock}{Teaching note}
Sinh viên nên xác định ba thành phần trong mỗi câu lệnh: đầu vào, phép toán và đầu ra.
\end{alertblock}
\end{frame}

\begin{frame}[fragile]{Ví dụ ngắn: giải thích từng dòng}
\small
\begin{lstlisting}[style=py]
import numpy as np

scores = np.array([7, 8, 6, 9])
bonus = scores + 1
average = scores.mean()
print(bonus)
print(average)
\end{lstlisting}
\begin{itemize}
  \item \texttt{import numpy as np}: nạp NumPy và đặt tên ngắn là \texttt{np}.
  \item \texttt{np.array([...])}: chuyển list Python thành mảng NumPy.
  \item \texttt{scores + 1}: cộng 1 vào mọi phần tử bằng vector hóa.
  \item \texttt{scores.mean()}: tính giá trị trung bình của mảng.
  \item \texttt{print(...)}: hiển thị kết quả.
\end{itemize}
\end{frame}

\section{NumPy là gì?}

\begin{frame}{NumPy là gì?}
\small
\key{NumPy}, viết tắt của \key{Numerical Python}, là thư viện Python cốt lõi cho tính toán số nhanh.
\vspace{0.4em}

Đối tượng trung tâm của nó là \key{N-dimensional array}, called \texttt{ndarray}. Một mảng NumPy có thể biểu diễn:
\begin{itemize}
  \item Một vectơ một chiều.
  \item Một ma trận hai chiều.
  \item Một tensor ba chiều.
  \item Các cấu trúc số có số chiều cao hơn.
\end{itemize}
\begin{block}{Vì sao quan trọng}
NumPy lưu trữ dữ liệu đồng nhất hiệu quả và hỗ trợ phép toán nhanh trên toàn bộ mảng.
\end{block}
\end{frame}

\begin{frame}{NumPy cung cấp gì?}
\small
\begin{columns}[T]
\column{0.47\linewidth}
\begin{itemize}
  \item Các phép toán mảng nhanh.
  \item Tính toán vector hóa.
  \item Broadcasting.
  \item Các hàm đại số tuyến tính.
\end{itemize}
\column{0.47\linewidth}
\begin{itemize}
  \item Các hàm thống kê.
  \item Sinh số ngẫu nhiên.
  \item Tích hợp với Pandas, SciPy, Matplotlib và Scikit-learn.
\end{itemize}
\end{columns}
\vspace{0.7em}
% \imageplaceholder{numpy\_linear\_algebra.webp}{NumPy linear algebra overview}
\end{frame}

\begin{frame}{Các đặc trưng chính}
\scriptsize
\begin{tabular}{p{0.26\linewidth}p{0.66\linewidth}}
\toprule
\textbf{Đặc trưng} & \textbf{Ý nghĩa} \\\\
\midrule
\texttt{ndarray} & Mảng nhiều chiều hiệu quả với một kiểu dữ liệu chung. \\\\
Vectorization & Áp dụng phép tính lên toàn bộ mảng mà không cần vòng lặp Python tường minh. \\\\
Broadcasting & Thực hiện phép toán trên các mảng có shape khác nhau nhưng tương thích. \\\\
Linear algebra & Nhân ma trận, định thức, nghịch đảo, trị riêng và tích vectơ. \\\\
Statistics & Trung bình, trung vị, phương sai, độ lệch chuẩn, percentile và quantile. \\\\
Integration & Làm việc với Pandas, SciPy, Matplotlib, Scikit-learn và Statsmodels. \\\\
\bottomrule
\end{tabular}
\end{frame}

\begin{frame}{Kiểm tra nhanh: NumPy là gì?}
\small
\quiztitle{Câu 1.} Cấu trúc dữ liệu chính trong NumPy là gì?
\begin{multicols}{2}
\begin{itemize}
  \item A. \texttt{DataFrame}
  \item B. \texttt{ndarray}
  \item C. \texttt{dictionary}
  \item D. \texttt{tuple}
\end{itemize}
\end{multicols}
\vspace{0.8em}
\quiztitle{Câu 2. Đúng hay sai?} NumPy được thiết kế chủ yếu cho tính toán số.
\end{frame}

\section{Tại sao nên học NumPy?}

\begin{frame}{Tại sao nên học NumPy?}
\small
NumPy cung cấp nền tảng hiệu quả cho công việc tính toán số trong Python.
\begin{itemize}
  \item Thực hiện phép toán vector hóa nhanh hơn nhiều vòng lặp Python thông thường.
  \item Lưu trữ dữ liệu số đồng nhất gọn hơn list Python chuẩn.
  \item Cung cấp hàm tối ưu cho ma trận, đại số tuyến tính và biến đổi.
  \item Hỗ trợ sinh số ngẫu nhiên và phân tích thống kê.
  \item Biểu diễn phép toán toán học phức tạp bằng cú pháp ngắn gọn.
  \item Là nền tảng tính toán số cho nhiều thư viện khoa học dữ liệu.
\end{itemize}
\end{frame}

\begin{frame}[fragile]{List Python và mảng NumPy}
\begin{columns}[T]
\column{0.48\linewidth}
\textbf{List Python: các đối tượng hỗn hợp}
\begin{lstlisting}[style=py]
values = [10, 2.5, "Python", True]
\end{lstlisting}
\begin{itemize}
  \item Cấu trúc chứa đối tượng linh hoạt.
  \item Không tối ưu cho mảng số lớn.
\end{itemize}
\column{0.48\linewidth}
\textbf{NumPy array: các giá trị đồng nhất}
\begin{lstlisting}[style=py]
import numpy as np
values = np.array([10, 20, 30, 40])
\end{lstlisting}
\begin{itemize}
  \item Cấu trúc đều đặn.
  \item Tính toán số hiệu quả.
\end{itemize}
\end{columns}
\end{frame}

\begin{frame}[fragile]{Ví dụ: nhân các giá trị}
\begin{columns}[T]
\column{0.48\linewidth}
\textbf{List Python}
\begin{lstlisting}[style=py]
values = [1, 2, 3, 4]
result = []

for value in values:
    result.append(value * 10)

print(result)
# [10, 20, 30, 40]
\end{lstlisting}
\column{0.48\linewidth}
\textbf{Vector hóa NumPy}
\begin{lstlisting}[style=py]
import numpy as np
values = np.array([1, 2, 3, 4])

result = values * 10
print(result)
# [10 20 30 40]
\end{lstlisting}
\end{columns}
\begin{block}{Ý tưởng chính}
Cùng một phép toán được áp dụng cho tất cả phần tử cùng lúc.
\end{block}
\end{frame}

\begin{frame}{Kiểm tra nhanh: tại sao học NumPy?}
\small
\quiztitle{Câu 1.} Vì sao mảng NumPy hiệu quả cho tính toán số?
\begin{itemize}
  \item A. Chúng luôn chứa văn bản
  \item B. Chúng dùng cấu trúc mảng đồng nhất và được tối ưu
  \item C. Chúng không dùng bộ nhớ
  \item D. Chúng tự động kết nối internet
\end{itemize}
\vspace{0.6em}
\quiztitle{Câu 2.} Biểu thức nào nhân mọi phần tử của mảng NumPy \texttt{a} by 10?
\begin{multicols}{2}
\begin{itemize}
  \item A. \texttt{a * 10}
  \item B. \texttt{a.append(10)}
  \item C. \texttt{a.add("10")}
  \item D. \texttt{a.sort(10)}
\end{itemize}
\end{multicols}
\end{frame}

\section{Tạo mảng}

\begin{frame}[fragile]{Tạo mảng NumPy}
\small
Mảng NumPy có thể được tạo từ list, tuple, range của Python hoặc các hàm dựng sẵn của NumPy.
\begin{lstlisting}[style=py]
import numpy as np

a = [9, 3, 3, 5]
arr = np.array(a)
print(arr)
# [9 3 3 5]
\end{lstlisting}
\vspace{0.2em}
\begin{columns}[T]
\column{0.5\linewidth}
\begin{lstlisting}[style=py]
arr = np.array([10, 20, 30, 40])
print(arr)
\end{lstlisting}
\column{0.5\linewidth}
\begin{lstlisting}[style=py]
matrix = np.array([
    [1, 2, 3],
    [4, 5, 6]
])
print(matrix)
\end{lstlisting}
\end{columns}
\end{frame}

\begin{frame}[fragile]{Mảng nhiều chiều}
\small
Mảng ba chiều có thể biểu diễn các ma trận xếp chồng hoặc tensor.
\begin{lstlisting}[style=py]
tensor = np.array([
    [
        [1, 2],
        [3, 4]
    ],
    [
        [5, 6],
        [7, 8]
    ]
])
print(tensor)
\end{lstlisting}
\begin{block}{Diễn giải}
Các chiều thường tương ứng với quan sát, đặc trưng, thời gian, kênh hoặc kịch bản.
\end{block}
\end{frame}

\begin{frame}[fragile]{Các hàm tạo mảng thường dùng}
\small
\begin{columns}[T]
\column{0.48\linewidth}
\begin{lstlisting}[style=py]
np.zeros(5)
np.zeros((2, 3))
np.ones((2, 3))
np.full((2, 3), 7)
\end{lstlisting}
\column{0.48\linewidth}
\begin{lstlisting}[style=py]
np.arange(0, 10, 2)
# [0 2 4 6 8]

np.linspace(0, 1, 5)
# [0.   0.25 0.5  0.75 1.]

np.eye(3)
\end{lstlisting}
\end{columns}
\end{frame}

\begin{frame}{Các lệnh chính: tạo mảng}
\small
\begin{tabular}{p{0.28\linewidth}p{0.61\linewidth}}
\toprule
\textbf{Lệnh} & \textbf{Ý nghĩa và khi nào dùng} \\
\midrule
\texttt{np.array(data)} & Chuyển list, tuple hoặc chuỗi lồng nhau thành \texttt{ndarray}. \\
\texttt{np.zeros(shape)} & Tạo mảng toàn số 0; hữu ích khi khởi tạo. \\
\texttt{np.ones(shape)} & Tạo mảng toàn số 1; hữu ích cho mask, trọng số hoặc khởi tạo. \\
\texttt{np.full(shape, value)} & Tạo mảng được lấp đầy bằng một giá trị cho trước. \\
\texttt{np.arange(start, stop, step)} & Sinh giá trị cách đều; \texttt{stop} không được lấy. \\
\texttt{np.linspace(start, stop, n)} & Sinh đúng \texttt{n} giá trị cách đều, mặc định bao gồm hai đầu mút. \\
\texttt{np.eye(n)} & Tạo ma trận đơn vị \(n\times n\). \\
\bottomrule
\end{tabular}
\end{frame}

\begin{frame}{Kiểm tra nhanh: tạo mảng}
\small
\quiztitle{Câu 1.} Hàm nào chuyển list Python thành mảng NumPy?
\begin{multicols}{2}
\begin{itemize}
  \item A. \texttt{np.array()}
  \item B. \texttt{np.list()}
  \item C. \texttt{np.convert()}
  \item D. \texttt{np.ndarray\_list()}
\end{itemize}
\end{multicols}
\vspace{0.4em}
\quiztitle{Câu 2.} Hàm nào tạo các giá trị cách đều giữa hai đầu mút?
\begin{multicols}{2}
\begin{itemize}
  \item A. \texttt{np.linspace()}
  \item B. \texttt{np.stack()}
  \item C. \texttt{np.split()}
  \item D. \texttt{np.mean()}
\end{itemize}
\end{multicols}
\end{frame}

\section{Thuộc tính mảng}

\begin{frame}[fragile]{Thuộc tính mảng: quan sát một mảng}
\small
\begin{lstlisting}[style=py]
arr = np.array([
    [1, 2, 3],
    [4, 5, 6]
])

print(arr.ndim)     # 2
print(arr.shape)    # (2, 3)
print(arr.size)     # 6
print(arr.dtype)
print(arr.itemsize)
\end{lstlisting}
\begin{block}{How to read these commands}
Mỗi câu lệnh yêu cầu NumPy mô tả mảng: số chiều, shape, tổng số phần tử, kiểu dữ liệu và bộ nhớ mỗi phần tử.
\end{block}
\end{frame}

\begin{frame}{Thuộc tính mảng: ý nghĩa từng thuộc tính}
\small
\begin{tabular}{p{0.22\linewidth}p{0.68\linewidth}}
\toprule
\textbf{Thuộc tính} & \textbf{Ý nghĩa} \\
\midrule
\texttt{ndim} & Số chiều. Vectơ có 1 chiều; ma trận có 2 chiều. \\
\texttt{shape} & Kích thước theo từng chiều. Ví dụ, \texttt{(2, 3)} nghĩa là 2 hàng và 3 cột. \\
\texttt{size} & Tổng số phần tử trong mảng. \\
\texttt{dtype} & Kiểu dữ liệu lưu trong mảng, chẳng hạn số nguyên hoặc số thực. \\
\texttt{itemsize} & Số byte dùng cho mỗi phần tử. \\
\bottomrule
\end{tabular}
\end{frame}

\begin{frame}{Các lệnh chính: kiểm tra thuộc tính mảng}
\small
\begin{tabular}{p{0.25\linewidth}p{0.64\linewidth}}
\toprule
\textbf{Thuộc tính} & \textbf{Nó cho biết điều gì} \\
\midrule
\texttt{arr.ndim} & Số trục/số chiều của mảng. \\
\texttt{arr.shape} & Độ dài của mảng theo từng trục; ví dụ, \texttt{(3,4)} means 3 rows and 4 columns. \\
\texttt{arr.size} & Tổng số phần tử được lưu trữ. \\
\texttt{arr.dtype} & Kiểu dùng để lưu mỗi phần tử, chẳng hạn \texttt{int64} or \texttt{float64}. \\
\texttt{arr.itemsize} & Number of bytes used by one element. \\
\texttt{arr.nbytes} & Tổng số byte dùng bởi tất cả phần tử mảng. \\
\bottomrule
\end{tabular}
\begin{block}{Thói quen hữu ích}
Trước khi reshape, broadcasting hoặc mô hình hóa, hãy kiểm tra \texttt{shape} and \texttt{dtype}; nhiều lỗi NumPy bắt nguồn từ shape hoặc kiểu dữ liệu không như mong đợi.
\end{block}
\end{frame}

\begin{frame}{Kiểm tra nhanh: thuộc tính mảng}
\small
\quiztitle{Câu 1.} Thuộc tính nào trả về kích thước của mảng?
\begin{multicols}{2}
\begin{itemize}
  \item A. \texttt{shape}
  \item B. \texttt{mean}
  \item C. \texttt{append}
  \item D. \texttt{index}
\end{itemize}
\end{multicols}
\vspace{0.4em}
\quiztitle{Câu 2.} Thuộc tính nào trả về tổng số phần tử?
\begin{multicols}{2}
\begin{itemize}
  \item A. \texttt{size}
  \item B. \texttt{dtype}
  \item C. \texttt{ndim}
  \item D. \texttt{itemsize}
\end{itemize}
\end{multicols}
\end{frame}

\section{Indexing và slicing}

\begin{frame}[fragile]{Indexing mảng}
\begin{columns}[T]
\column{0.48\linewidth}
\textbf{Indexing một chiều}
\begin{lstlisting}[style=py]
arr = np.array([10, 20, 30, 40])

print(arr[0])   # 10
print(arr[2])   # 30
print(arr[-1])  # 40
\end{lstlisting}
\column{0.48\linewidth}
\textbf{Indexing hai chiều}
\begin{lstlisting}[style=py]
matrix = np.array([
    [1, 2, 3],
    [4, 5, 6]
])

print(matrix[0, 1])  # 2
print(matrix[1, 2])  # 6
\end{lstlisting}
\end{columns}
\begin{block}{Quy ước indexing}
Python và NumPy dùng chỉ số bắt đầu từ 0.
\end{block}
\end{frame}

\begin{frame}[fragile]{Slicing mảng}
\begin{columns}[T]
\column{0.48\linewidth}
\textbf{Slicing một chiều}
\begin{lstlisting}[style=py]
arr = np.array([10, 20, 30, 40, 50])

print(arr[1:4])
# [20 30 40]

print(arr[::2])
# [10 30 50]
\end{lstlisting}
\column{0.48\linewidth}
\textbf{Slicing hai chiều}
\begin{lstlisting}[style=py]
matrix = np.array([
    [1, 2, 3],
    [4, 5, 6],
    [7, 8, 9]
])

print(matrix[0:2, 1:3])
# [[2 3]
#  [5 6]]
\end{lstlisting}
\end{columns}
\end{frame}

\begin{frame}[fragile]{View và bản sao}
\small
Nhiều phép slicing của NumPy trả về \key{view} của mảng gốc thay vì một bản sao độc lập.
\begin{lstlisting}[style=py]
arr = np.array([10, 20, 30, 40])

view = arr[1:3]
view[0] = 999

print(arr)
# [ 10 999  30  40]
\end{lstlisting}
Tạo bản sao độc lập khi cần:
\begin{lstlisting}[style=py]
copy = arr[1:3].copy()
\end{lstlisting}
\begin{alertblock}{Quan trọng}
Thay đổi một view có thể làm thay đổi mảng gốc.
\end{alertblock}
\end{frame}

\begin{frame}{Các lệnh chính: indexing và slicing}
\small
\begin{tabular}{p{0.28\linewidth}p{0.61\linewidth}}
\toprule
\textbf{Cú pháp} & \textbf{Ý nghĩa} \\
\midrule
\texttt{arr[i]} & Chọn phần tử ở vị trí \texttt{i}; chỉ số âm đếm từ cuối mảng. \\
\texttt{arr[a:b]} & Chọn các vị trí \texttt{a} through \texttt{b-1}. \\
\texttt{arr[a:b:s]} & Slice with step size \texttt{s}. \\
\texttt{matrix[i,j]} & Select row \texttt{i}, column \texttt{j}. \\
\texttt{matrix[r1:r2,c1:c2]} & Trích xuất một mảng con hình chữ nhật. \\
\texttt{arr[mask]} & Boolean indexing: retain values where \texttt{mask} is \texttt{True}. \\
\texttt{arr.copy()} & Make independent data when a view should not modify the original. \\
\bottomrule
\end{tabular}
\end{frame}

\begin{frame}{Kiểm tra nhanh: indexing và slicing}
\small
\quiztitle{Câu 1.} Điều gì diễn ra khi \texttt{arr[0]} trả về gì?
\begin{multicols}{2}
\begin{itemize}
  \item A. The first element
  \item B. The last element
  \item C. The array shape
  \item D. Kích thước mảng
\end{itemize}
\end{multicols}
\quiztitle{Câu 2.} In a two-dimensional array, \texttt{matrix[1, 2]} refers to:
\begin{itemize}
  \item A. Row index 1 and column index 2
  \item B. Row 1 only
  \item C. Column 2 only
  \item D. The array dimensions
\end{itemize}
\quiztitle{Câu 3. Đúng hay sai?} Một lát cắt NumPy có thể dùng chung bộ nhớ với mảng gốc.
\end{frame}

\section{Reshaping, stacking, splitting}

\begin{frame}[fragile]{Đổi hình dạng mảng}
\small
Reshape thay đổi kích thước/chiều của mảng mà không thay đổi dữ liệu.
\begin{lstlisting}[style=py]
arr = np.arange(12)
matrix = arr.reshape(3, 4)
print(matrix)
# [[ 0  1  2  3]
#  [ 4  5  6  7]
#  [ 8  9 10 11]]
\end{lstlisting}
Dùng \texttt{-1} để suy ra một chiều:
\begin{lstlisting}[style=py]
matrix = arr.reshape(2, -1)
\end{lstlisting}
Làm phẳng mảng:
\begin{lstlisting}[style=py]
flat = ma trận.flatten()
flat_view = ma trận.ravel()
\end{lstlisting}
\smallnote{\texttt{flatten()} trả về bản sao; \texttt{ravel()} thường trả về view nếu có thể.}
\end{frame}

\begin{frame}[fragile]{Thay đổi kích thước mảng}
\small
\texttt{resize()} thay đổi shape và có thể thay đổi số phần tử.
\begin{lstlisting}[style=py]
arr = np.array([1, 2, 3, 4])
resized = np.resize(arr, (2, 3))
print(resized)
# [[1 2 3]
#  [4 1 2]]
\end{lstlisting}
\begin{block}{So sánh}
\begin{itemize}
  \item \texttt{reshape()} yêu cầu tổng số phần tử phải tương thích.
  \item \texttt{resize()} có thể tạo tổng kích thước khác bằng cách lặp lại giá trị.
\end{itemize}
\end{block}
\end{frame}

\begin{frame}[fragile]{Ghép mảng}
\begin{columns}[T]
\column{0.5\linewidth}
\begin{lstlisting}[style=py]
a = np.array([1, 2, 3])
b = np.array([4, 5, 6])

np.vstack((a, b))
# [[1 2 3]
#  [4 5 6]]
\end{lstlisting}
\column{0.5\linewidth}
\begin{lstlisting}[style=py]
np.hstack((a, b))
# [1 2 3 4 5 6]

np.stack((a, b), axis=0)
\end{lstlisting}
\end{columns}
\begin{block}{Mục đích}
Stacking ghép các mảng thành mảng lớn hơn theo một trục được chọn.
\end{block}
\end{frame}

\begin{frame}[fragile]{Tách mảng}
\small
Splitting chia mảng thành các mảng nhỏ hơn.
\begin{lstlisting}[style=py]
arr = np.array([1, 2, 3, 4, 5, 6])
parts = np.split(arr, 3)
print(parts)
# [array([1, 2]), array([3, 4]), array([5, 6])]
\end{lstlisting}
Đối với mảng nhiều chiều:
\begin{lstlisting}[style=py]
matrix = np.array([
    [1, 2, 3, 4],
    [5, 6, 7, 8]
])

left, right = np.hsplit(matrix, 2)
\end{lstlisting}
\end{frame}

\begin{frame}{Các lệnh chính: reshape, stack và split}
\scriptsize
\begin{tabular}{p{0.29\linewidth}p{0.60\linewidth}}
\toprule
\textbf{Lệnh} & \textbf{Ý nghĩa} \\
\midrule
\texttt{arr.reshape(r,c)} & Sắp xếp lại cùng các phần tử vào shape mới tương thích. \\
\texttt{arr.reshape(...,-1)} & Yêu cầu NumPy tự suy ra một chiều. \\
\texttt{arr.flatten()} & Trả về dạng phẳng 1-D \emph{copy}. \\
\texttt{arr.ravel()} & Trả về view phẳng khi có thể. \\
\texttt{np.resize(arr, shape)} & Thay đổi tổng kích thước nếu cần; giá trị có thể lặp lại. \\
\texttt{np.vstack((a,b))} & Ghép mảng theo chiều dọc (hướng hàng). \\
\texttt{np.hstack((a,b))} & Ghép mảng theo chiều ngang. \\
\texttt{np.stack((a,b), axis=k)} & Thêm trục mới và ghép theo trục đó. \\
\texttt{np.split(arr,n)} & Chia mảng thành \texttt{n} phần bằng nhau khi có thể. \\
\texttt{np.hsplit(matrix,n)} & Chia ma trận theo chiều cột. \\
\bottomrule
\end{tabular}
\end{frame}

\begin{frame}{Kiểm tra nhanh: reshaping và splitting}
\small
\quiztitle{Câu 1.} Phương thức nào đổi mảng từ shape này sang shape khác?
\begin{multicols}{2}
\begin{itemize}
  \item A. \texttt{reshape()}
  \item B. \texttt{mean()}
  \item C. \texttt{split()}
  \item D. \texttt{sort\_index()}
\end{itemize}
\end{multicols}
\quiztitle{Câu 2.} Điều gì diễn ra khi \texttt{-1} nghĩa là gì trong \texttt{reshape()}?
\begin{itemize}
  \item A. NumPy nên tự suy ra chiều đó
  \item B. Xóa một chiều
  \item C. Đảo ngược mảng
  \item D. Chuyển giá trị thành số âm
\end{itemize}
\quiztitle{Câu 3.} Hàm nào ghép mảng theo chiều dọc? \texttt{np.vstack()}, \texttt{np.mean()}, \texttt{np.split()}, or \texttt{np.random()}?
\end{frame}

\section{Broadcasting và số học}

\begin{frame}[fragile]{Broadcasting: ý tưởng cốt lõi}
\small
Broadcasting cho phép NumPy thực hiện phép toán giữa các mảng có shape tương thích.
\begin{lstlisting}[style=py]
arr = np.array([1, 2, 3, 4])
result = arr + 10
print(result)
# [11 12 13 14]
\end{lstlisting}
\begin{block}{Quy tắc cơ bản}
Bắt đầu từ các chiều ngoài cùng bên phải, hai chiều tương thích khi:
\begin{itemize}
  \item chúng bằng nhau; hoặc
  \item một trong chúng là \texttt{1}.
\end{itemize}
\end{block}
\end{frame}

\begin{frame}[fragile]{Broadcasting giữa các mảng}
\small
\begin{lstlisting}[style=py]
matrix = np.array([
    [1, 2, 3],
    [4, 5, 6]
])

row = np.array([10, 20, 30])
result = matrix + row
print(result)
# [[11 22 33]
#  [14 25 36]]
\end{lstlisting}
\begin{alertblock}{Shape không tương thích}
\begin{lstlisting}[style=py]
a = np.ones((2, 3))
b = np.ones((2, 2))
# a + b gây lỗi broadcasting
\end{lstlisting}
\end{alertblock}
\end{frame}

\begin{frame}[fragile]{Các phép toán số học cơ bản}
\small
NumPy hỗ trợ số học theo từng phần tử.
\begin{lstlisting}[style=py]
a = np.array([1, 2, 3])
b = np.array([4, 5, 6])

print(a + b)   # [5 7 9]
print(a - b)   # [-3 -3 -3]
print(a * b)   # [ 4 10 18]
print(a / b)
print(a ** 2)  # [1 4 9]
print(b % a)
\end{lstlisting}
\begin{block}{Phân biệt quan trọng}
\texttt{\*} thực hiện nhân theo từng phần tử, không phải nhân ma trận.
\end{block}
\end{frame}

\begin{frame}{Các lệnh chính: số học và broadcasting}
\small
\begin{tabular}{p{0.28\linewidth}p{0.61\linewidth}}
\toprule
\textbf{Biểu thức} & \textbf{Ý nghĩa} \\
\midrule
\texttt{a + b}, \texttt{a - b} & Cộng và trừ theo từng phần tử. \\
\texttt{a * b} & Nhân theo từng phần tử. \\
\texttt{a / b} & Chia theo từng phần tử. \\
\texttt{a ** p} & Nâng mọi phần tử lên lũy thừa \texttt{p}. \\
\texttt{a \% b} & Phần dư theo từng phần tử. \\
\texttt{array + scalar} & Broadcast số vô hướng đến mọi phần tử. \\
\texttt{matrix + row} & Broadcast một hàng 1-D tương thích qua các hàng của ma trận. \\
\bottomrule
\end{tabular}
\begin{alertblock}{Kiểm tra shape trước}
Broadcasting so sánh các chiều từ phải sang trái. Các chiều phải bằng nhau hoặc một trong chúng phải bằng 1.
\end{alertblock}
\end{frame}

\begin{frame}{Kiểm tra nhanh: broadcasting và số học}
\small
\quiztitle{Câu 1.} Broadcasting cho phép điều gì?
\begin{itemize}
  \item A. Phép toán trên các mảng có shape khác nhau nhưng tương thích
  \item B. Truyền dữ liệu tự động qua internet
  \item C. Chuyển mảng thành tệp
  \item D. Loại bỏ mọi chiều
\end{itemize}
\quiztitle{Câu 2.} Hai chiều tương thích broadcasting khi chúng bằng nhau hoặc:
\begin{multicols}{2}
\begin{itemize}
  \item A. Một trong chúng bằng 1
  \item B. Cả hai đều âm
  \item C. Cả hai là văn bản
  \item D. Tổng của chúng bằng 0
\end{itemize}
\end{multicols}
\quiztitle{Câu 3.} Điều gì diễn ra khi \texttt{a * b} làm gì với các mảng cùng shape?
\end{frame}

\section{Tổng hợp và universal functions}

\begin{frame}[fragile]{Các hàm tổng hợp}
\small
Các hàm tổng hợp tóm tắt một mảng.
\begin{lstlisting}[style=py]
arr = np.array([9, 3, 3, 5])

print(arr.sum())   # 20
print(arr.mean())  # 5.0
print(arr.min())
print(arr.max())
print(np.median(arr))
print(np.var(arr))
print(np.std(arr))
\end{lstlisting}
\begin{block}{Cách dùng điển hình}
Aggregation biến mảng thô thành thống kê tóm tắt.
\end{block}
\end{frame}

\begin{frame}[fragile]{Tổng hợp theo trục}
\small
\begin{lstlisting}[style=py]
matrix = np.array([
    [1, 2, 3],
    [4, 5, 6]
])

print(ma trận.sum(axis=0))
# [5 7 9]

print(ma trận.sum(axis=1))
# [ 6 15]
\end{lstlisting}
\begin{itemize}
  \item \texttt{axis=0}: tổng hợp theo chiều hàng, tạo kết quả theo cột.
  \item \texttt{axis=1}: tổng hợp theo chiều cột, tạo kết quả theo hàng.
\end{itemize}
\end{frame}

\begin{frame}[fragile]{Universal functions}
\small
Universal function, hay \key{ufunc}, thực hiện phép toán theo từng phần tử trên mảng.
\begin{columns}[T]
\column{0.48\linewidth}
\begin{lstlisting}[style=py]
arr = np.array([1, 4, 9, 16])
print(np.sqrt(arr))
# [1. 2. 3. 4.]

print(np.exp(np.array([0, 1, 2])))
\end{lstlisting}
\column{0.48\linewidth}
\begin{lstlisting}[style=py]
values = np.array([1, np.e, np.e**2])
print(np.log(values))

angles = np.array([0, np.pi / 2, np.pi])
print(np.sin(angles))
\end{lstlisting}
\end{columns}
\vspace{0.4em}
Các ví dụ khác: \texttt{np.abs()}, \texttt{np.round()}.
\end{frame}

\begin{frame}{Các lệnh chính: tổng hợp và universal functions}
\scriptsize
\begin{tabular}{p{0.29\linewidth}p{0.60\linewidth}}
\toprule
\textbf{Lệnh} & \textbf{Ý nghĩa} \\
\midrule
\texttt{arr.sum()}, \texttt{arr.mean()} & Tổng và trung bình số học. \\
\texttt{arr.min()}, \texttt{arr.max()} & Giá trị nhỏ nhất và lớn nhất. \\
\texttt{np.median(arr)} & Trung vị. \\
\texttt{np.var(arr)}, \texttt{np.std(arr)} & Phương sai và độ lệch chuẩn. \\
\texttt{func(..., axis=0)} & Tổng hợp theo chiều hàng, tạo một kết quả cho mỗi cột. \\
\texttt{func(..., axis=1)} & Tổng hợp theo chiều cột, tạo một kết quả cho mỗi hàng. \\
\texttt{np.sqrt(arr)} & Căn bậc hai theo từng phần tử. \\
\texttt{np.exp(arr)}, \texttt{np.log(arr)} & Hàm mũ và log tự nhiên theo từng phần tử. \\
\texttt{np.sin(arr)} & Hàm sin theo từng phần tử; góc tính bằng radian. \\
\texttt{np.abs(arr)}, \texttt{np.round(arr)} & Giá trị tuyệt đối và làm tròn. \\
\bottomrule
\end{tabular}
\end{frame}

\begin{frame}{Kiểm tra nhanh: tổng hợp và ufunc}
\small
\quiztitle{Câu 1.} Hàm nào tính trung bình số học?
\begin{multicols}{2}
\begin{itemize}
  \item A. \texttt{np.mean()}
  \item B. \texttt{np.stack()}
  \item C. \texttt{np.dot()}
  \item D. \texttt{np.reshape()}
\end{itemize}
\end{multicols}
\quiztitle{Câu 2.} Trong mảng hai chiều, \texttt{axis=0} thường tổng hợp theo chiều nào?
\begin{itemize}
  \item A. Theo hàng, tạo kết quả theo cột
  \item B. Theo cột, tạo kết quả theo hàng
  \item C. Tên tệp
  \item D. Kiểu dữ liệu
\end{itemize}
\quiztitle{Câu 3.} Hàm nào tính căn bậc hai theo từng phần tử?
\end{frame}

\section{Toán tử Boolean}

\begin{frame}[fragile]{So sánh và lọc Boolean}
\small
NumPy hỗ trợ so sánh theo từng phần tử.
\begin{lstlisting}[style=py]
arr = np.array([10, 20, 30, 40])
print(arr > 20)
# [False False  True  True]

selected = arr[arr > 20]
print(selected)
# [30 40]
\end{lstlisting}
Kết hợp điều kiện:
\begin{lstlisting}[style=py]
selected = arr[(arr >= 20) & (arr <= 30)]
print(selected)
# [20 30]
\end{lstlisting}
Dùng \texttt{\&} cho AND, \texttt{|} cho OR và \texttt{\textasciitilde} cho NOT.
\end{frame}

\begin{frame}{Các lệnh chính: toán tử Boolean}
\small
\begin{tabular}{p{0.31\linewidth}p{0.58\linewidth}}
\toprule
\textbf{Biểu thức} & \textbf{Ý nghĩa} \\
\midrule
\texttt{arr > value} & Tạo mảng Boolean từ phép so sánh theo từng phần tử. \\
\texttt{arr[arr > value]} & Lọc và chỉ trả về các giá trị thỏa điều kiện. \\
\texttt{(cond1) \& (cond2)} & AND logic theo từng phần tử. \\
\texttt{(cond1) | (cond2)} & OR logic theo từng phần tử. \\
\texttt{\textasciitilde mask} & NOT logic theo từng phần tử. \\
\texttt{np.where(cond,a,b)} & Chọn \texttt{a} nơi điều kiện đúng và \texttt{b} trong trường hợp ngược lại. \\
\bottomrule
\end{tabular}
\begin{block}{Cú pháp quan trọng}
Dùng dấu ngoặc quanh các phép so sánh kết hợp vì \texttt{\&} and \texttt{|} có độ ưu tiên khác với toán tử so sánh.
\end{block}
\end{frame}

\begin{frame}{Kiểm tra nhanh: toán tử Boolean}
\small
\quiztitle{Câu 1.} Điều gì diễn ra khi \texttt{arr[arr > 20]} trả về gì?
\begin{itemize}
  \item A. Các phần tử lớn hơn 20
  \item B. Kích thước mảng
  \item C. Kiểu dữ liệu của mảng
  \item D. Tất cả phần tử được chuyển thành giá trị Boolean
\end{itemize}
\vspace{0.8em}
\quiztitle{Thảo luận.} Vì sao lọc Boolean hữu ích trong phân tích dữ liệu?
\end{frame}

\section{Đại số tuyến tính}

\begin{frame}[fragile]{Nhân ma trận}
\small
\begin{lstlisting}[style=py]
A = np.array([
    [1, 2],
    [3, 4]
])

result = np.dot(A, A)
print(result)
# [[ 7 10]
#  [15 22]]

result = A @ A
\end{lstlisting}
\begin{block}{Nhân theo từng phần tử và nhân ma trận}
\begin{itemize}
  \item \texttt{A \* A}: nhân theo từng phần tử.
  \item \texttt{A @ A}: nhân ma trận.
\end{itemize}
\end{block}
\end{frame}

\begin{frame}[fragile]{Các phép toán ma trận}
\small
\begin{columns}[T]
\column{0.48\linewidth}
\begin{lstlisting}[style=py]
print(A.T)

determinant = np.linalg.det(A)
print(determinant)

inverse = np.linalg.inv(A)
print(inverse)
\end{lstlisting}
\column{0.48\linewidth}
\begin{lstlisting}[style=py]
A = np.array([
    [2, 1],
    [1, 3]
])

b = np.array([8, 13])
x = np.linalg.giải(A, b)
print(x)
\end{lstlisting}
\end{columns}
\begin{block}{Điều kiện}
Nghịch đảo chỉ tồn tại với ma trận vuông không suy biến.
\end{block}
\end{frame}

\begin{frame}[fragile]{Trị riêng và tích vectơ}
\small
\begin{columns}[T]
\column{0.48\linewidth}
\begin{lstlisting}[style=py]
eigenvalues, eigenvectors = np.linalg.eig(A)

print(eigenvalues)
print(eigenvectors)
\end{lstlisting}
\column{0.48\linewidth}
\begin{lstlisting}[style=py]
a = np.array([1, 2, 3])
b = np.array([4, 5, 6])

print(np.inner(a, b))
print(np.outer(a, b))
print(np.dot(a, b))
print(np.vdot(a, b))
\end{lstlisting}
\end{columns}
\smallnote{Với mảng một chiều thực, \texttt{dot} and \texttt{vdot} thường giống nhau. \texttt{vdot} cũng xử lý liên hợp phức.}
\end{frame}

\begin{frame}{Các lệnh chính: đại số tuyến tính}
\scriptsize
\begin{tabular}{p{0.29\linewidth}p{0.60\linewidth}}
\toprule
\textbf{Lệnh} & \textbf{Ý nghĩa} \\
\midrule
\texttt{A @ B} & Nhân ma trận; cú pháp dễ đọc được ưu tiên. \\
\texttt{np.dot(A,B)} & Tích vô hướng / tích ma trận tùy theo số chiều. \\
\texttt{A.T} & Chuyển vị mảng. \\
\texttt{np.linalg.det(A)} & Tính định thức của ma trận vuông. \\
\texttt{np.linalg.inv(A)} & Tính nghịch đảo của ma trận vuông không suy biến. \\
\texttt{np.linalg.giải(A,b)} & Solve \(Ax=b\) trực tiếp; thường tốt hơn việc tính \(A^{-1}b\). \\
\texttt{np.linalg.eig(A)} & Trả về trị riêng và vectơ riêng. \\
\texttt{np.inner(a,b)} & Tích trong. \\
\texttt{np.outer(a,b)} & Tích ngoài. \\
\texttt{np.vdot(a,b)} & Tích vô hướng vectơ có liên hợp phức khi cần. \\
\bottomrule
\end{tabular}
\end{frame}

\begin{frame}{Kiểm tra nhanh: đại số tuyến tính}
\small
\quiztitle{Câu 1.} Toán tử nào thực hiện nhân ma trận?
\begin{multicols}{2}
\begin{itemize}
  \item A. \texttt{@}
  \item B. \texttt{\%}
  \item C. \texttt{//}
  \item D. \texttt{\&}
\end{itemize}
\end{multicols}
\quiztitle{Câu 2.} Hàm nào tính nghịch đảo ma trận?
\begin{itemize}
  \item A. \texttt{np.linalg.inv()}
  \item B. \texttt{np.mean()}
  \item C. \texttt{np.resize()}
  \item D. \texttt{np.random()}
\end{itemize}
\quiztitle{Câu 3.} Giải thích sự khác nhau giữa \texttt{A * B} and \texttt{A @ B}.
\end{frame}

\section{Số ngẫu nhiên và thống kê}

\begin{frame}[fragile]{Sinh số ngẫu nhiên}
\small
Với mã mới, NumPy thường dùng một đối tượng sinh số ngẫu nhiên.
\begin{lstlisting}[style=py]
rng = np.random.default_rng()
\end{lstlisting}
\begin{columns}[T]
\column{0.48\linewidth}
\begin{lstlisting}[style=py]
rng = np.random.default_rng(42)
values = rng.integers(
    low=1,
    high=10,
    size=5
)
print(values)
\end{lstlisting}
\column{0.48\linewidth}
\begin{lstlisting}[style=py]
values = rng.normal(
    loc=0,
    scale=1,
    size=5
)
print(values)
\end{lstlisting}
\end{columns}
\begin{alertblock}{Lưu ý}
Bộ sinh số ngẫu nhiên của NumPy dùng cho tính toán số và thống kê; không nên xem là bộ sinh bảo mật mật mã.
\end{alertblock}
\end{frame}

\begin{frame}[fragile]{Các phân phối thường dùng}
\small
\begin{lstlisting}[style=py]
rng = np.random.default_rng(42)

rng.uniform(low=0, high=1, size=5)
rng.normal(loc=0, scale=1, size=5)
rng.binomial(n=10, p=0.5, size=5)
rng.poisson(lam=3, size=5)
rng.exponential(scale=2, size=5)
rng.chisquare(df=4, size=5)
\end{lstlisting}
\begin{block}{Tính tái lập}
Dùng cùng seed giúp tái tạo cùng chuỗi giả ngẫu nhiên trong môi trường tương thích.
\end{block}
\end{frame}

\begin{frame}[fragile]{Các hàm thống kê}
\small
\begin{lstlisting}[style=py]
data = np.array([12, 15, 18, 20, 25])

print(np.mean(data))
print(np.median(data))
print(np.var(data))
print(np.std(data))
print(np.percentile(data, 25))
print(np.quantile(data, 0.75))
\end{lstlisting}
\begin{block}{Quy ước tổng thể và mẫu}
Mặc định, \texttt{np.var()} and \texttt{np.std()} use \texttt{ddof=0}. Với quy ước mẫu thường dùng, hãy dùng \texttt{ddof=1}.
\end{block}
\end{frame}

\begin{frame}[fragile]{Giá trị số bị thiếu}
\small
NumPy thường biểu diễn giá trị số bị thiếu bằng \texttt{np.nan}.
\begin{lstlisting}[style=py]
data = np.array([10.0, 20.0, np.nan, 40.0])
print(np.mean(data))
# nan

print(np.nanmean(data))
print(np.nanmedian(data))
print(np.nanstd(data))

print(np.isnan(data))
clean_data = data[~np.isnan(data)]
\end{lstlisting}
\smallnote{Các quy trình xử lý dữ liệu thiếu đầy đủ hơn thường được thực hiện bằng Pandas.}
\end{frame}

\begin{frame}{Các lệnh chính: số ngẫu nhiên và thống kê}
\scriptsize
\begin{tabular}{p{0.32\linewidth}p{0.57\linewidth}}
\toprule
\textbf{Lệnh} & \textbf{Ý nghĩa} \\
\midrule
\texttt{np.random.default\_rng(seed)} & Tạo bộ sinh số ngẫu nhiên hiện đại; seed hỗ trợ tái lập. \\
\texttt{rng.integers(low,high,size)} & Sinh số nguyên ngẫu nhiên; \texttt{high} không được lấy. \\
\texttt{rng.uniform(...)} & Lấy mẫu từ phân phối đều. \\
\texttt{rng.normal(loc,scale,size)} & Lấy mẫu từ phân phối chuẩn với trung bình \texttt{loc} and độ lệch chuẩn \texttt{scale}. \\
\texttt{rng.binomial(...)} / \texttt{poisson(...)} & Lấy mẫu từ các phân phối rời rạc phổ biến. \\
\texttt{np.percentile(data,q)} & Trả về percentile \texttt{q} trên thang 0--100. \\
\texttt{np.quantile(data,q)} & Trả về quantile \texttt{q} trên thang 0--1. \\
\texttt{np.std(data,ddof=1)} & Quy ước độ lệch chuẩn mẫu. \\
\texttt{np.isnan(data)} & Xác định \texttt{NaN} giá trị. \\
\texttt{np.nanmean(data)} & Trung bình khi bỏ qua \texttt{NaN}. \\
\bottomrule
\end{tabular}
\end{frame}

\begin{frame}{Kiểm tra nhanh: số ngẫu nhiên và thống kê}
\small
\quiztitle{Câu 1.} Phương thức nào tạo bộ sinh số ngẫu nhiên hiện đại của NumPy?
\begin{itemize}
  \item A. \texttt{np.random.default\_rng()}
  \item B. \texttt{np.random.file()}
  \item C. \texttt{np.create\_random\_array()}
  \item D. \texttt{np.random.text()}
\end{itemize}
\quiztitle{Câu 2.} Vì sao seed hữu ích?
\begin{itemize}
  \item A. Nó giúp tái tạo chuỗi giả ngẫu nhiên
  \item B. Nó làm cho giá trị thực sự không thể dự đoán
  \item C. Nó loại bỏ mọi giá trị ngẫu nhiên
  \item D. Nó chuyển dữ liệu thành chuỗi
\end{itemize}
\quiztitle{Câu 3.} Tham số nào cho quy ước độ lệch chuẩn mẫu thường dùng? \texttt{ddof=1}, \texttt{axis=-100}, \texttt{dtype="sample"}, or \texttt{copy=False}?
\end{frame}

\section{Vector hóa và hiệu năng}

\begin{frame}[fragile]{Các phép toán vector hóa}
\small
Vectorization áp dụng phép toán lên toàn bộ mảng cùng lúc.
\begin{columns}[T]
\column{0.48\linewidth}
\begin{lstlisting}[style=py]
a = np.arange(5)
result = a * 10
print(result)
# [ 0 10 20 30 40]
\end{lstlisting}
\column{0.48\linewidth}
\begin{lstlisting}[style=py]
a = list(range(5))
result = []
for value in a:
    result.append(value * 10)
print(result)
\end{lstlisting}
\end{columns}
\begin{block}{Lợi ích điển hình}
Mã ngắn hơn, dễ đọc hơn, nhanh hơn với mảng số lớn và phù hợp hơn với quy trình khoa học.
\end{block}
\end{frame}

\begin{frame}[fragile]{Logic điều kiện vector hóa}
\small
\begin{lstlisting}[style=py]
arr = np.array([10, 20, 30, 40])

labels = np.where(
    arr >= 30,
    "high",
    "low"
)

print(labels)
# ['low' 'low' 'high' 'high']
\end{lstlisting}
\begin{block}{Diễn giải}
Điều kiện vector hóa hữu ích cho gán nhãn, sàng lọc, lọc và biến đổi dữ liệu.
\end{block}
\end{frame}

\begin{frame}{Các lệnh chính: vector hóa}
\small
\begin{tabular}{p{0.30\linewidth}p{0.59\linewidth}}
\toprule
\textbf{Mẫu lệnh} & \textbf{Ý nghĩa} \\
\midrule
\texttt{arr * scalar} & Áp dụng phép toán số học cho mọi phần tử mà không cần vòng lặp Python tường minh. \\
\texttt{np.sqrt(arr)} & Áp dụng ufunc theo từng phần tử trong mã NumPy đã tối ưu. \\
\texttt{arr >= threshold} & Đánh giá điều kiện cho toàn mảng cùng lúc. \\
\texttt{np.where(cond,a,b)} & Biến đổi điều kiện theo kiểu vector hóa. \\
\bottomrule
\end{tabular}
\begin{block}{Ý tưởng cốt lõi}
Vectorization mô tả \emph{cách ta biểu diễn tính toán}: phép toán được viết cho mảng thay vì lặp thủ công qua từng phần tử.
\end{block}
\end{frame}

\begin{frame}{Kiểm tra nhanh: vector hóa}
\small
\quiztitle{Câu 1.} Ưu điểm chính của vector hóa là gì?
\begin{itemize}
  \item A. Nó áp dụng phép toán mảng mà không cần vòng lặp Python tường minh
  \item B. Nó chuyển mọi mảng thành văn bản
  \item C. Nó loại bỏ nhu cầu dùng bộ nhớ
  \item D. Nó ngăn mọi lỗi
\end{itemize}
\vspace{0.8em}
\quiztitle{Thảo luận.} Khi nào mã vector hóa dễ bảo trì hơn mã dựa trên vòng lặp?
\end{frame}

\section{Bộ nhớ, sắp xếp, ảnh}

\begin{frame}[fragile]{Lưu ý về bộ nhớ}
\small
Mảng NumPy thường lưu giá trị đồng nhất trong bố cục bộ nhớ đều đặn.
\begin{lstlisting}[style=py]
arr = np.array([1, 2, 3, 4], dtype=np.int32)
print(arr.nbytes)

small_values = np.array([1, 2, 3, 4], dtype=np.int8)
print(small_giá trị.dtype)
print(small_giá trị.nbytes)
\end{lstlisting}
\begin{block}{Lựa chọn kiểu dữ liệu}
Kiểu dữ liệu nhỏ hơn có thể giảm bộ nhớ, nhưng phải biểu diễn được miền giá trị và độ chính xác cần thiết.
\end{block}
\end{frame}

\begin{frame}[fragile]{Chuyển đổi kiểu dữ liệu}
\small
\begin{lstlisting}[style=py]
arr = np.array([1.2, 2.8, 3.5])
integers = arr.astype(np.int32)
print(integers)
# [1 2 3]
\end{lstlisting}
\begin{alertblock}{Quan trọng}
Chuyển số thực sang số nguyên sẽ loại bỏ phần thập phân.
\end{alertblock}
\end{frame}

\begin{frame}[fragile]{Sắp xếp và tìm kiếm}
\small
\begin{lstlisting}[style=py]
arr = np.array([9, 3, 7, 1])
sorted_arr = np.sort(arr)
print(sorted_arr)
# [1 3 7 9]

arr = np.array([10, 20, 30, 40])
indices = np.where(arr > 20)
print(indices)

arr = np.array([1, 2, 2, 3, 3, 3])
values, counts = np.unique(arr, return_counts=True)
print(values)
print(counts)
\end{lstlisting}
\end{frame}

\begin{frame}[fragile]{Làm việc với ảnh}
\small
Ảnh số có thể được biểu diễn bằng mảng NumPy.
\begin{itemize}
  \item Ảnh xám: mảng hai chiều.
  \item Ảnh màu: mảng ba chiều chứa chiều cao, chiều rộng và kênh màu.
\end{itemize}
\begin{lstlisting}[style=py]
image = np.array([
    [0, 128, 255],
    [255, 128, 0],
    [50, 100, 150]
])
print(image.shape)
# (3, 3)

brighter = np.clip(image + 30, 0, 255)
\end{lstlisting}
\smallnote{\texttt{np.clip()} giữ các giá trị trong khoảng hợp lệ.}
\end{frame}

\begin{frame}{Các lệnh chính: bộ nhớ, kiểu dữ liệu, sắp xếp và tìm kiếm}
\scriptsize
\begin{tabular}{p{0.30\linewidth}p{0.59\linewidth}}
\toprule
\textbf{Lệnh} & \textbf{Ý nghĩa} \\
\midrule
\texttt{arr.nbytes} & Tổng bộ nhớ dùng bởi phần tử mảng. \\
\texttt{arr.astype(dtype)} & Chuyển giá trị mảng sang kiểu dữ liệu mới. \\
\texttt{np.sort(arr)} & Trả về giá trị đã sắp xếp. \\
\texttt{np.where(condition)} & Trả về chỉ số thỏa điều kiện khi chỉ cung cấp điều kiện. \\
\texttt{np.unique(arr)} & Trả về giá trị duy nhất. \\
\texttt{np.unique(arr, return\_counts=True)} & Trả về giá trị duy nhất và tần suất. \\
\texttt{np.clip(arr,low,high)} & Giới hạn mọi giá trị trong khoảng chỉ định; hữu ích cho cường độ ảnh. \\
\bottomrule
\end{tabular}
\end{frame}

\begin{frame}{Kiểm tra nhanh: bộ nhớ, tìm kiếm, ảnh}
\small
\quiztitle{Câu 1.} Thuộc tính nào trả về số byte dùng bởi phần tử mảng?
\begin{multicols}{2}
\begin{itemize}
  \item A. \texttt{nbytes}
  \item B. \texttt{ndim}
  \item C. \texttt{mean}
  \item D. \texttt{shape}
\end{itemize}
\end{multicols}
\quiztitle{Câu 2.} Hàm nào trả về các giá trị duy nhất của mảng?
\begin{itemize}
  \item A. \texttt{np.unique()}
  \item B. \texttt{np.reshape()}
  \item C. \texttt{np.random()}
  \item D. \texttt{np.outer()}
\end{itemize}
\quiztitle{Câu 3.} Ảnh xám thường được biểu diễn bằng loại mảng nào?
\end{frame}

\section{Tích hợp và lỗi}

\begin{frame}[fragile]{Tích hợp với Pandas}
\small
Một Pandas \texttt{Series} or \texttt{DataFrame} có thể trao đổi dữ liệu với NumPy.
\begin{lstlisting}[style=py]
import numpy as np
import pandas as pd

arr = np.array([
    [1, 2],
    [3, 4],
    [5, 6]
])

df = pd.DataFrame(arr, columns=["A", "B"])
array_from_df = df.to_numpy()

df["A_sqrt"] = np.sqrt(df["A"])
\end{lstlisting}
\end{frame}

\begin{frame}[fragile]{Tích hợp với Scikit-learn}
\small
Nhiều mô hình Scikit-learn chấp nhận mảng NumPy làm đầu vào.
\begin{lstlisting}[style=py]
from sklearn.linear_model import LinearRegression

X = np.array([[1], [2], [3], [4]])
y = np.array([2, 4, 6, 8])

model = LinearRegression()
model.fit(X, y)

prediction = model.predict(np.array([[5]]))
print(prediction)
\end{lstlisting}
NumPy hỗ trợ ma trận đặc trưng, vectơ mục tiêu, dự đoán mô hình, tiền xử lý số và phép tính đánh giá.
\end{frame}

\begin{frame}[fragile]{Các lỗi NumPy thường gặp}
\small
\begin{columns}[T]
\column{0.48\linewidth}
\textbf{Shape không khớp}
\begin{lstlisting}[style=py]
a = np.ones((2, 3))
b = np.ones((2, 2))
# a + b raises an error
\end{lstlisting}
\textbf{Reshape không hợp lệ}
\begin{lstlisting}[style=py]
arr = np.arange(10)
# arr.reshape(3, 4) raises an error
\end{lstlisting}
\column{0.48\linewidth}
\textbf{Chia cho 0}
\begin{lstlisting}[style=py]
arr = np.array([1.0, 0.0])
print(1 / arr)
\end{lstlisting}
\textbf{Tràn số nguyên}
\begin{lstlisting}[style=py]
arr = np.array([127], dtype=np.int8)
print(arr + 1)
\end{lstlisting}
\end{columns}
\end{frame}

\begin{frame}{Thực hành tốt}
\small
\begin{itemize}
  \item Import NumPy bằng \texttt{import numpy as np}.
  \item Dùng phép toán vector hóa thay cho vòng lặp Python khi phù hợp.
  \item Check \texttt{shape}, \texttt{ndim}, and \texttt{dtype} trước các phép toán phức tạp.
  \item Chỉ dùng broadcasting khi hành vi shape mong muốn đã rõ.
  \item Phân biệt nhân theo phần tử với nhân ma trận.
  \item Dùng \texttt{.copy()} khi cần mảng độc lập.
  \item Chọn kiểu dữ liệu cẩn thận để cân bằng miền giá trị, độ chính xác và bộ nhớ.
  \item Dùng seed cố định khi tính tái lập quan trọng.
  \item Treat \texttt{NaN}, vô cực và tràn số một cách tường minh.
\end{itemize}
\end{frame}

\begin{frame}{Các lệnh chính: tích hợp với công cụ khoa học dữ liệu}
\small
\begin{tabular}{p{0.32\linewidth}p{0.57\linewidth}}
\toprule
\textbf{Lệnh} & \textbf{Vai trò} \\
\midrule
\texttt{pd.DataFrame(arr,...)} & Tạo DataFrame Pandas từ mảng NumPy. \\
\texttt{df.to\_numpy()} & Trích xuất giá trị DataFrame thành mảng NumPy. \\
\texttt{np.sqrt(df["A"])} & Áp dụng ufunc NumPy cho dữ liệu Pandas tương thích. \\
\texttt{model.fit(X,y)} & Huấn luyện mô hình Scikit-learn bằng dữ liệu đặc trưng và mục tiêu dạng mảng. \\
\texttt{model.predict(Xnew)} & Tạo dự đoán cho quan sát mới dạng mảng. \\
\bottomrule
\end{tabular}
\begin{block}{Mối liên hệ}
Mảng NumPy là định dạng trao đổi dữ liệu số phổ biến giữa Pandas, SciPy, Matplotlib và Scikit-learn.
\end{block}
\end{frame}

\begin{frame}{Kiểm tra nhanh: tích hợp và lỗi}
\small
\quiztitle{Câu 1.} Phương thức Pandas nào chuyển DataFrame thành mảng NumPy?
\begin{multicols}{2}
\begin{itemize}
  \item A. \texttt{to\_numpy()}
  \item B. \texttt{to\_list\_only()}
  \item C. \texttt{as\_matrix\_text()}
  \item D. \texttt{convert\_numpy\_file()}
\end{itemize}
\end{multicols}
\quiztitle{Câu 2.} Trong học máy, mảng NumPy hai chiều thường biểu diễn:
\begin{itemize}
  \item A. Ma trận đặc trưng
  \item B. Tên tệp
  \item C. Tiêu đề biểu đồ
  \item D. Trình cài đặt package
\end{itemize}
\quiztitle{Câu 3.} Vì sao \texttt{np.arange(10).reshape(3, 4)} bị lỗi?
\end{frame}

\section{Tóm tắt và thực hành}

\begin{frame}{Tóm tắt nội dung I}
\small
\begin{tabular}{p{0.25\linewidth}p{0.66\linewidth}}
\toprule
\textbf{Chủ đề} & \textbf{Ý chính} \\
\midrule
NumPy & Thư viện Python cốt lõi cho tính toán số. \\
\texttt{ndarray} & Mảng đồng nhất N chiều. \\
Vectorization & Phép toán mảng không cần vòng lặp Python tường minh. \\
Broadcasting & Phép toán trên các shape khác nhau nhưng tương thích. \\
Indexing và slicing & Truy cập và trích xuất phần tử mảng. \\
Reshaping & Thay đổi số chiều/kích thước mảng. \\
\bottomrule
\end{tabular}
\end{frame}

\begin{frame}{Tóm tắt nội dung II}
\small
\begin{tabular}{p{0.25\linewidth}p{0.66\linewidth}}
\toprule
\textbf{Chủ đề} & \textbf{Ý chính} \\
\midrule
Aggregation & Tổng, trung bình, nhỏ nhất, lớn nhất, phương sai và độ lệch chuẩn. \\
Universal functions & Các hàm toán học nhanh theo từng phần tử. \\
Linear algebra & Nhân ma trận, nghịch đảo, định thức và trị riêng. \\
Random generation & Sinh giá trị từ các phân phối xác suất. \\
Statistics & Trung bình, trung vị, phương sai, độ lệch chuẩn và percentile. \\
Integration & Làm việc chặt chẽ với Pandas, SciPy và Scikit-learn. \\
\bottomrule
\end{tabular}
\end{frame}

\begin{frame}{Thực hành: tạo mảng}
\small
\begin{block}{Bài tập 1}
\begin{enumerate}
  \item Tạo mảng chứa các giá trị từ 1 đến 20.
  \item Reshape nó thành \(4 \times 5\) ma trận.
  \item In shape, số chiều, size và kiểu dữ liệu.
\end{enumerate}
\end{block}
\begin{block}{Bài tập 2}
Sử dụng \(4 \times 5\) matrix:
\begin{enumerate}
  \item Trích xuất hàng đầu tiên.
  \item Trích xuất cột cuối cùng.
  \item Trích xuất phần trung tâm \(2 \times 3\) .
  \item Chọn tất cả giá trị chẵn bằng lọc Boolean.
\end{enumerate}
\end{block}
\end{frame}

\begin{frame}{Thực hành: thống kê và broadcasting}
\small
\begin{block}{Bài tập 3: Statistics}
Tạo mảng gồm mười giá trị số và tính trung bình, trung vị, phương sai, độ lệch chuẩn mẫu, nhỏ nhất, lớn nhất và các percentile 25, 50, 75.
\end{block}
\begin{block}{Bài tập 4: Broadcasting}
\begin{enumerate}
  \item Tạo một ma trận \(3 \times 4\).
  \item Tạo một mảng một chiều với bốn giá trị.
  \item Cộng mảng một chiều vào từng hàng của ma trận.
  \item Giải thích vì sao broadcasting hợp lệ.
\end{enumerate}
\end{block}
\end{frame}

\begin{frame}{Thực hành: đại số tuyến tính và giá trị thiếu}
\small
\begin{block}{Bài tập 5: Linear algebra}
Given
\[
A=\begin{bmatrix}2&1\\ 1&3\end{bmatrix},\quad b=\begin{bmatrix}8\\ 13\end{bmatrix},
\]
tính định thức của \(A\), nghịch đảo của \(A\), giải \(Ax=b\), và kiểm tra bằng \texttt{A @ x}.
\end{block}
\begin{block}{Bài tập 6: Missing values}
Given \texttt{[10.0, np.nan, 20.0, 30.0, np.nan, 40.0]}, đếm giá trị thiếu, tính trung bình khi bỏ qua giá trị thiếu, loại bỏ giá trị thiếu và thay thế giá trị thiếu bằng trung vị.
\end{block}
\end{frame}

\begin{frame}{Câu hỏi suy ngẫm cuối bài}
\small
\begin{itemize}
  \item Chủ đề NumPy nào quan trọng nhất cho quy trình phân tích dữ liệu?
  \item Khi nào nên kiểm tra \texttt{shape} trước khi chạy phép toán?
  \item Vectorization và broadcasting giảm vòng lặp thủ công như thế nào?
  \item Vì sao cần phân biệt view và bản sao?
  \item Khi nào nên dùng \texttt{ddof=1} thay cho mặc định?
  \item Lỗi NumPy thường gặp nào bạn dễ gặp nhất trong thực hành?
\end{itemize}
\end{frame}

\end{document}